\chapter{Dedekind Domains}

\begin{center}\itshape
In the integers of a number field, numbers may factor in several ways, but ideals factor in exactly one. This chapter finds the rings where that happens, measures by a group how far they are from having
unique factorisation of numbers, and shows that numbers and curves obey the same algebra.
\end{center}

In 1871 Richard Dedekind added a supplement, the tenth, to the second edition of Dirichlet's \emph{Vorlesungen über Zahlentheorie}, the lectures of his late friend that he had edited. Dirichlet had
died in 1859, and the book was by then as much Dedekind's as Dirichlet's. In the supplement Dedekind replaced Kummer's mysterious ``ideal numbers'' by sets of ordinary numbers, which he called \emph{ideals}, and proved that every
ideal of the ring of integers of a number field is a product of prime ideals in exactly one way. Volume~13 tells how this rescued arithmetic after the unique factorisation of numbers had failed, and proves the theorem for
rings of algebraic integers, using the fact that their quotients are finite.

Dedekind was a quiet man who spent almost his whole career at the technical school of his home town, Braunschweig, and never married; he lived with his sister. He also lived long. According to E.~T.~Bell, a
\emph{Calendar for mathematicians} once gave the date of his death as 4 September 1899, and Dedekind wrote to the editor that he had spent that day in perfect health, in stimulating conversation with his guest and
friend Georg Cantor. He actually died in 1916, aged 84.

Two later steps made his theory into the one in this chapter. In 1882 Dedekind and Heinrich Weber applied the same ideal theory to algebraic \emph{functions}, that is to curves: the ring $\C[x,y]/(y^2-x^3+x)$ of
polynomial functions on a cubic curve factors its ideals exactly like a ring of algebraic integers. And in 1927 Emmy Noether found the axioms. She showed that unique factorisation of ideals holds precisely in the
rings that are Noetherian, integrally closed and of dimension one, three properties that Chapters~4 and~5 have prepared. Chapter~5 decomposes every proper ideal of a Noetherian ring into primary ideals. In a Dedekind domain that decomposition is the prime factorisation. When colleagues praised her for this kind of work, she would answer, according to a story
that Leo Corry calls ``mathematical lore'', \emph{Es steht alles schon bei Dedekind}: it is all already in Dedekind.

All rings are commutative with $1$. Throughout, $R$ is a domain with field of fractions $K$.

\section{Fractional ideals}

The rational number $\frac12$ is the inverse of $2$; to invert ideals, one needs ideals with denominators.

\begin{definition}\label{def:fractional}
A \emph{fractional ideal} of $R$ is a non-zero $R$-submodule $I\subseteq K$ such that $dI\subseteq R$ for some non-zero $d\in R$. The product $IJ$ is the submodule generated by the products $xy$, $x\in I$, $y\in J$,
and
\[
(R:I)=\{x\in K: xI\subseteq R\}.
\]
$I$ is \emph{invertible} if $IJ=R$ for some fractional ideal $J$.
\end{definition}

Non-zero ideals of $R$ are fractional ideals (with $d=1$), and so is $xR$ for every $x\in K^\times$: these are the \emph{principal} fractional ideals, and $xR\cdot x^{-1}R=R$. If $R$ is Noetherian, every fractional
ideal is finitely generated, since $I\cong dI$ is an ideal. The product is associative and commutative with unit $R$. In $\Z$, the fractional ideals are the $\frac abZ$, and $\frac ab\Z\cdot\frac ba\Z=\Z$.

\begin{lemma}\label{lem:invertible}
\leavevmode
\begin{enumerate}[label=(\alph*)]
\item If $IJ=R$, then $J=(R:I)$, and $I$ is finitely generated. One writes $I^{-1}=(R:I)$.
\item A product $IJ$ is invertible if and only if $I$ and $J$ are.
\item If $R$ is local, a fractional ideal is invertible if and only if it is principal.
\item If $I$ is finitely generated, then $I$ is invertible if and only if $I_{\mathfrak m}$ is invertible in $R_{\mathfrak m}$ for every maximal ideal $\mathfrak m$.
\end{enumerate}
\end{lemma}

\begin{proof}
(a) Clearly $J\subseteq(R:I)$, and $(R:I)=(R:I)IJ\subseteq RJ=J$, because $(R:I)I\subseteq R$. Write $1=\sum_{i=1}^na_ib_i$ with $a_i\in I$, $b_i\in J$. For $x\in I$, $x=\sum(xb_i)a_i$ with $xb_i\in IJ=R$, so
$a_1,\dots,a_n$ generate $I$.

(b) If $IJL=R$, then $I$ has inverse $JL$ and $J$ has inverse $IL$. Conversely $IJ\cdot I^{-1}J^{-1}=R$.

(c) Principal fractional ideals are invertible. Conversely, with $1=\sum a_ib_i$ as in (a), the products $a_ib_i$ lie in $R$, and they cannot all lie in the maximal ideal, since their sum is $1$. So some $u=a_ib_i$ is a unit,
and for $x\in I$, $x=(xb_i)u^{-1}a_i\in a_iR$. Hence $I=a_iR$.

(d) $I(R:I)$ is an ideal of $R$, and $I$ is invertible if and only if it equals $R$, that is, if and only if it lies in no maximal ideal, that is (Proposition~\ref{prop:localzero} applied to $R/I(R:I)$), if and only if
$\bigl(I(R:I)\bigr)_{\mathfrak m}=I_{\mathfrak m}(R:I)_{\mathfrak m}$ equals $R_{\mathfrak m}$ for every $\mathfrak m$. It remains to see that $(R:I)_{\mathfrak m}=(R_{\mathfrak m}:I_{\mathfrak m})$ when $I=(a_1,\dots,a_n)$.
The inclusion $\subseteq$ is clear. If $xI_{\mathfrak m}\subseteq R_{\mathfrak m}$, then $xa_i=r_i/s_i$ with $s_i\notin\mathfrak m$, and $s=s_1\cdots s_n$ satisfies $sxa_i\in R$ for all $i$, so $sx\in(R:I)$ and
$x=sx/s\in(R:I)_{\mathfrak m}$. Here the finite generation of $I$ was essential.
\end{proof}

\section{Discrete valuation rings}

Near a single prime, a Dedekind domain looks like $\Z_{(p)}$: every non-zero element is a unit times a power of $p$.

\begin{definition}\label{def:dvr}
A \emph{discrete valuation ring} (DVR) is a principal ideal domain with exactly one non-zero prime ideal.
\end{definition}

If $R$ is a DVR with maximal ideal $\mathfrak m=(\pi)$, then $\pi$ is, up to units, the only irreducible element, so by unique factorisation (Volume~3, Theorem~6.14) every $x\in K^\times$ is $x=u\pi^n$ with $u\in R^\times$ and
$n\in\Z$ unique. The map $v(x)=n$ is a homomorphism $K^\times\to\Z$ with $v(x+y)\ge\min(v(x),v(y))$, a \emph{discrete valuation}, and $R=\{0\}\cup\{x:v(x)\ge0\}$. The ideals of $R$ are $0$ and the $\mathfrak m^n=(\pi^n)$.
Examples: $\Z_{(p)}$ with the $p$-adic valuation; $\C[x]_{(x-a)}$ with the order of vanishing at $a$; the formal power series $k[[t]]$ (Exercise~\ref{ex:powerseries}).

\begin{theorem}\label{thm:dvr}
Let $(R,\mathfrak m)$ be a Noetherian local domain whose only non-zero prime ideal is $\mathfrak m\ne0$. The following are equivalent:
\begin{enumerate}[label=(\alph*)]
\item $R$ is a discrete valuation ring;
\item $R$ is integrally closed;
\item $\mathfrak m$ is principal;
\item $\mathfrak m$ is invertible.
\end{enumerate}
Moreover (c) implies (a) for every Noetherian local domain with $\mathfrak m\ne0$.
\end{theorem}

\begin{proof}
(a)$\Rightarrow$(b). Let $x=u\pi^n\in K$ be integral over $R$, with $x^k+c_{k-1}x^{k-1}+\dots+c_0=0$, $c_i\in R$. If $n<0$, multiply by $\pi^{-nk}$: then $u^k=-\sum_{i<k}c_iu^i\pi^{-n(k-i)}\in\mathfrak m$, which is absurd since
$u^k$ is a unit. So $n\ge0$ and $x\in R$.

(c)$\Leftrightarrow$(d) is Lemma~\ref{lem:invertible}(c).

(c)$\Rightarrow$(a), using only that $R$ is a Noetherian local domain and $\mathfrak m=(\pi)\ne0$. First, $\bigcap_n\mathfrak m^n=0$. Indeed, if $x=\pi^ny_n$ for every $n$, then $\pi y_{n+1}=y_n$, so
$(y_1)\subseteq(y_2)\subseteq\cdots$, and this chain is stationary (Proposition~\ref{prop:noetherian}): $y_{n+1}=cy_n=c\pi y_{n+1}$ for some $n$ and $c\in R$. As $1-c\pi\notin\mathfrak m$ is a unit, $y_{n+1}=0$ and $x=0$. So every
$x\ne0$ lies in a largest power $\mathfrak m^n$, $x=\pi^nu$ with $u\notin\mathfrak m$, a unit. A non-zero ideal $I$ is then $(\pi^n)$, where $n$ is the least exponent occurring in $I$. So $R$ is a principal ideal domain, and
its non-zero primes $(\pi^n)$ must have $n=1$ (as $\pi\cdot\pi^{n-1}\in(\pi^n)$).

(b)$\Rightarrow$(c). Choose $x\in\mathfrak m$, $x\ne0$. The ring $R/(x)$ is Noetherian and its only prime ideal is $\mathfrak m/(x)$, since the primes of $R$ containing $x$ are non-zero. Its nilradical, the intersection of its primes
(Exercise~\ref{ex:localiff}), is therefore $\mathfrak m/(x)$, a finitely generated ideal made of nilpotent elements. If $\bar a_1,\dots,\bar a_k$ generate it, with $\bar a_i^{n_i}=0$, then any product of $n_1+\dots+n_k$ of
its elements is $0$, by expanding and pigeonholing. Hence $\mathfrak m^N\subseteq(x)$ for some $N\ge1$, and we choose $N$ minimal. Pick $y\in\mathfrak m^{N-1}\setminus(x)$ (with $\mathfrak m^0=R$) and put $z=y/x\in K$. Then
$z\notin R$, as $y\notin(x)$, while $z\mathfrak m\subseteq R$, as $y\mathfrak m\subseteq\mathfrak m^N\subseteq(x)$.

Suppose $z\mathfrak m\subseteq\mathfrak m$. Write $\mathfrak m=(m_1,\dots,m_k)$ and $zm_i=\sum_ja_{ij}m_j$ with $a_{ij}\in R$. The matrix $zI-A$ kills the non-zero vector $(m_1,\dots,m_k)$ of $K^k$, so its determinant
vanishes, and expanding it gives a monic equation for $z$ over $R$: $z$ is integral over $R$, so $z\in R$ since $R$ is integrally closed, a contradiction. Hence $z\mathfrak m$ is an ideal of $R$ not contained in
$\mathfrak m$, so $z\mathfrak m=R$. Then $\mathfrak m=z^{-1}R$, and $z^{-1}=x/y\in\mathfrak m$: the maximal ideal is principal.
\end{proof}

\begin{think}
In (b)$\Rightarrow$(c), the element $z=y/x$ is ``almost'' in $R$: it fails to be integral only because it has a pole of order exactly one along $\mathfrak m$. In $\Z_{(p)}$ one may take $x=p^2$, $y=p$, $z=1/p$. Where in the proof
was the dimension hypothesis used, and why is it needed? (Think of $\C[x,y]_{(x,y)}$, integrally closed but with a maximal ideal that is not principal.)
\end{think}

\section{Dedekind domains}

\begin{definition}\label{def:dedekind}
A \emph{Dedekind domain} is a Noetherian, integrally closed domain, not a field, in which every non-zero prime ideal is maximal.
\end{definition}

Principal ideal domains that are not fields are Dedekind (they are UFDs, hence integrally closed by the argument of Theorem~\ref{thm:dvr}, (a)$\Rightarrow$(b); their non-zero primes $(p)$ are maximal). Volume~13, Chapter~4, shows
that the ring of integers $\mathcal O_K$ of a number field is Dedekind. Theorem~\ref{thm:intclosure} below gives both facts at once, and many more.

\begin{theorem}[Noether, 1927]\label{thm:dedekind}
For a Noetherian domain $R$ that is not a field, the following are equivalent:
\begin{enumerate}[label=(\alph*)]
\item $R$ is a Dedekind domain;
\item $R_{\mathfrak p}$ is a discrete valuation ring for every non-zero prime $\mathfrak p$;
\item every fractional ideal of $R$ is invertible;
\item every non-zero proper ideal of $R$ is a product of prime ideals.
\end{enumerate}
In that case the factorisation in (d) is unique up to the order of the factors.
\end{theorem}

\begin{proof}
(a)$\Rightarrow$(b). $R_{\mathfrak p}$ is a Noetherian local domain (an ideal $J$ of $R_{\mathfrak p}$ is generated by $J\cap R$, which is finitely generated). Its primes correspond to the primes of $R$ inside
$\mathfrak p$, which are $0$ and $\mathfrak p$. It is integrally closed: if $x\in K$ satisfies $x^k+\sum_{i<k}c_ix^i=0$ with $c_i\in R_{\mathfrak p}$, let $s\notin\mathfrak p$ be a common denominator of the $c_i$; multiplying by
$s^k$ shows that $sx$ is integral over $R$, so $sx\in R$ and $x\in R_{\mathfrak p}$. By Theorem~\ref{thm:dvr}, $R_{\mathfrak p}$ is a DVR.

(b)$\Rightarrow$(a). Every non-zero prime $\mathfrak p$ is maximal: if $\mathfrak p\subseteq\mathfrak m$, then $\mathfrak pR_{\mathfrak m}$ is a non-zero prime of the DVR $R_{\mathfrak m}$, so it is $\mathfrak mR_{\mathfrak m}$, and $\mathfrak p=\mathfrak m$
by the correspondence of Theorem~\ref{thm:locprimes}. Next, $R=\bigcap_{\mathfrak m}R_{\mathfrak m}$ inside $K$: if $x$ lies in every $R_{\mathfrak m}$, the ideal $\{r\in R:rx\in R\}$ is contained in no maximal ideal, so it contains $1$.
If $x\in K$ is integral over $R$, it is integral over each $R_{\mathfrak m}$, which is integrally closed, so $x\in\bigcap R_{\mathfrak m}=R$.

(b)$\Rightarrow$(c). A fractional ideal $I$ is finitely generated, and $I_{\mathfrak m}$ is a non-zero fractional ideal of the principal ideal domain $R_{\mathfrak m}$, hence principal and invertible. Apply Lemma~\ref{lem:invertible}(d).

(c)$\Rightarrow$(b). Let $\mathfrak p\ne0$ be prime. The ideal $\mathfrak p$ is invertible, $\mathfrak p\mathfrak p^{-1}=R$, so $\mathfrak p_{\mathfrak p}(\mathfrak p^{-1})_{\mathfrak p}=R_{\mathfrak p}$: the maximal ideal of $R_{\mathfrak p}$ is invertible,
hence principal (Lemma~\ref{lem:invertible}(c)), and $R_{\mathfrak p}$ is a DVR by the last sentence of Theorem~\ref{thm:dvr}.

(c)$\Rightarrow$(d). Suppose some non-zero proper ideal is not a product of primes. Since $R$ is Noetherian, there is a maximal such ideal $I$. It is not prime, so $I\subsetneq\mathfrak m$ for a maximal ideal $\mathfrak m$. Now
$I\subseteq I\mathfrak m^{-1}\subseteq\mathfrak m\mathfrak m^{-1}=R$. If $I\mathfrak m^{-1}=I$, multiplying by $I^{-1}$ gives $\mathfrak m^{-1}=R$, and then $R=\mathfrak m\mathfrak m^{-1}=\mathfrak m$, absurd. So $I\mathfrak m^{-1}$ is an
ideal strictly larger than $I$, and it is proper, since $I\mathfrak m^{-1}=R$ would give $I=\mathfrak m$. By maximality $I\mathfrak m^{-1}$ is a product of primes, and so is $I=(I\mathfrak m^{-1})\mathfrak m$, a contradiction.

(d)$\Rightarrow$(c). This is the delicate direction, and it goes in three steps.

\emph{Step 1: cancellation.} Suppose $P_1\cdots P_r=Q_1\cdots Q_s$ with all $P_i,Q_j$ prime and all $P_i$ invertible. Then $r=s$ and, after renumbering, $P_i=Q_i$. Induction on $r$: for $r=0$ the left side is $R$, so $s=0$.
Otherwise let $P_1$ be minimal among the $P_i$. As $Q_1\cdots Q_s\subseteq P_1$ and $P_1$ is prime, some $Q_j\subseteq P_1$, say $Q_1$. Likewise some $P_i\subseteq Q_1\subseteq P_1$, so $P_i=P_1$ by minimality, and $Q_1=P_1$. Multiply by $P_1^{-1}$ and use induction.

\emph{Step 2: an invertible prime $P$ is maximal.} Let $a\notin P$, and suppose $P+aR\ne R$. By (d), $P+aR=P_1\cdots P_r$ and $P+a^2R=Q_1\cdots Q_s$, with primes containing $P$. In the domain $\bar R=R/P$, where $\bar a\ne0$,
this reads $\bar a\bar R=\bar P_1\cdots\bar P_r$ and $\bar a^2\bar R=\bar Q_1\cdots\bar Q_s$, so $\bar Q_1\cdots\bar Q_s=\bar P_1^2\cdots\bar P_r^2$. The $\bar P_i$ are invertible by Lemma~\ref{lem:invertible}(b), being factors of the
principal ideal $\bar a\bar R$, so by Step~1 the $\bar Q_j$ are the $\bar P_i$, each taken twice. Since the primes of $R$ containing $P$ correspond to those of $\bar R$, the $Q_j$ are the $P_i$, each taken twice, and $P+a^2R=(P+aR)^2\subseteq P^2+aR$. So every $x\in P$ is $x=y+az$ with $y\in P^2$; then $az\in P$, so $z\in P$, and
$P\subseteq P^2+aP=P(P+aR)\subseteq P$. Multiplying $P=P(P+aR)$ by $P^{-1}$ gives $P+aR=R$, a contradiction.

\emph{Step 3.} Let $\mathfrak p\ne0$ be prime and $0\ne b\in\mathfrak p$. Write $bR=P_1\cdots P_r$. Each $P_i$ is invertible, as a factor of $bR$, hence maximal by Step~2. Since $P_1\cdots P_r\subseteq\mathfrak p$, some $P_i\subseteq\mathfrak p$,
and $P_i=\mathfrak p$ by maximality. So every non-zero prime is invertible, every non-zero ideal is invertible as a product of invertible primes, and every fractional ideal $I=d^{-1}(dI)$ is invertible.

Uniqueness follows from Step~1, since under (c) all primes are invertible.
\end{proof}

For a Dedekind domain, therefore, the fractional ideals form a group $\mathcal I(R)$ under multiplication, and it is the \emph{free abelian group on the non-zero primes}: every fractional ideal is uniquely
\[
I=\prod_{\mathfrak p}\mathfrak p^{v_{\mathfrak p}(I)},\qquad v_{\mathfrak p}(I)\in\Z,
\]
with finitely many non-zero exponents. (Write $I=d^{-1}(dI)$ and factor $dR$ and $dI$.) For $x\in K^\times$, $v_{\mathfrak p}(xR)$ is the valuation of $x$ in the DVR $R_{\mathfrak p}$, since $\mathfrak p$ is the only prime that survives in
$R_{\mathfrak p}$.

\begin{corollary}[to contain is to divide]\label{cor:contain}
Let $R$ be Dedekind and $I,J$ fractional ideals. Then $J\subseteq I$ if and only if $J=IL$ for an ideal $L\subseteq R$, if and only if $v_{\mathfrak p}(J)\ge v_{\mathfrak p}(I)$ for every $\mathfrak p$. Consequently
\[
v_{\mathfrak p}(I+J)=\min\bigl(v_{\mathfrak p}(I),v_{\mathfrak p}(J)\bigr),\qquad v_{\mathfrak p}(I\cap J)=\max\bigl(v_{\mathfrak p}(I),v_{\mathfrak p}(J)\bigr).
\]
\end{corollary}

\begin{proof}
If $J\subseteq I$, then $L=I^{-1}J\subseteq I^{-1}I=R$ and $J=IL$. If $J=IL$, then $v_{\mathfrak p}(J)=v_{\mathfrak p}(I)+v_{\mathfrak p}(L)\ge v_{\mathfrak p}(I)$, since an ideal has non-negative exponents. If the exponents compare, then $L=I^{-1}J$
has non-negative exponents, so $L\subseteq R$ and $J=IL\subseteq I$. Now $I+J$ is the smallest fractional ideal containing $I$ and $J$, so its exponents are the largest ones below both, the minima; likewise for $I\cap J$.
\end{proof}

Let $I=\prod_{i=1}^r\mathfrak p_i^{n_i}$ be a non-zero ideal, with distinct non-zero primes and $n_i\ge1$. Corollary~\ref{cor:contain} gives $I=\bigcap_i\mathfrak p_i^{n_i}$, because the two sides have the same valuation at every prime. Each $\mathfrak p_i^{n_i}$ is $\mathfrak p_i$-primary: its radical is $\mathfrak p_i$, since every element of $\mathfrak p_i$ has valuation at least $1$ and its $n_i$-th power has valuation at least $n_i$, while a power lying in $\mathfrak p_i^{n_i}$ lies in the prime $\mathfrak p_i$, and if $ab\in\mathfrak p_i^{n_i}$ with $a\notin\mathfrak p_i^{n_i}$, then $v_{\mathfrak p_i}(a)<n_i$, so $v_{\mathfrak p_i}(b)\ge1$ and $b^{n_i}\in\mathfrak p_i^{n_i}$. The valuation at $\mathfrak p_i$ of $\bigcap_{j\ne i}\mathfrak p_j^{n_j}$ is $0$, while the valuation of $\mathfrak p_i^{n_i}$ is $n_i$, so no factor contains the intersection of the others. Theorem~\ref{thm:assprimes} therefore says that the associated primes of $I$ are exactly $\mathfrak p_1,\dots,\mathfrak p_r$, the primes in its factorisation. Each of them is isolated. A non-zero prime is maximal, so no associated prime of a non-zero ideal contains another.

\begin{example}\label{ex:sixtwice}
In $R=\Z[\sqrt{-5}]$, which is Dedekind (Volume~13, Chapter~4), the two factorisations $6=2\cdot3=(1+\sqrt{-5})(1-\sqrt{-5})$ of the number $6$ come from one factorisation of the ideal $6R$. With
$\mathfrak p=(2,1+\sqrt{-5})$, $\mathfrak q=(3,1+\sqrt{-5})$ and $\bar{\mathfrak q}=(3,1-\sqrt{-5})$, one has (Exercise~\ref{ex:zsqrt5})
\[
2R=\mathfrak p^2,\quad 3R=\mathfrak q\bar{\mathfrak q},\quad (1+\sqrt{-5})=\mathfrak p\mathfrak q,\quad(1-\sqrt{-5})=\mathfrak p\bar{\mathfrak q},
\]
so both factorisations of $6$ regroup the same four prime ideals $\mathfrak p,\mathfrak p,\mathfrak q,\bar{\mathfrak q}$. None of these primes is principal.
\end{example}

\section{The class group}

\begin{definition}\label{def:classgroup}
For a Dedekind domain $R$, the \emph{class group} $\operatorname{Cl}(R)$ is the quotient of $\mathcal I(R)$ by the subgroup $\mathcal P(R)$ of principal fractional ideals. Its order, when finite, is the \emph{class number}.
\end{definition}

\begin{proposition}\label{prop:classseq}
The sequence of abelian groups
\[
1\longrightarrow R^\times\longrightarrow K^\times\xrightarrow{\ x\,\mapsto\, xR\ }\mathcal I(R)\longrightarrow\operatorname{Cl}(R)\longrightarrow1
\]
is exact, and $R$ is a principal ideal domain if and only if $\operatorname{Cl}(R)=1$, if and only if $R$ is a unique factorisation domain.
\end{proposition}

\begin{proof}
$xR=R$ if and only if $x\in R$ and $x^{-1}\in R$, so the kernel of $x\mapsto xR$ is $R^\times$; its image is $\mathcal P(R)$, the kernel of the projection. If $\operatorname{Cl}(R)=1$, every ideal is principal; conversely, a principal ideal domain
has trivial class group, since every fractional ideal is $d^{-1}(dI)$. A principal ideal domain is a UFD (Volume~3, Theorem~6.14). Finally, let $R$ be a UFD and $\mathfrak p\ne0$ a prime. A non-zero $x\in\mathfrak p$ is a product of
irreducibles, one of which, $\pi$, lies in $\mathfrak p$. In a UFD, $(\pi)$ is a non-zero prime ideal, hence maximal, so $\mathfrak p=(\pi)$. All primes being principal, so are all products of primes, that is all ideals.
\end{proof}

So $\operatorname{Cl}(R)$ measures exactly the failure of unique factorisation of elements. For $\Z[\sqrt{-5}]$ it has order $2$, generated by the class of $\mathfrak p$; Volume~13 computes such groups with Minkowski's geometry of
numbers and proves that they are finite for rings of integers. In general they need not be: Luther Claborn proved in 1966 that \emph{every} abelian group is the class group of some Dedekind domain. The next result shows that
even a non-principal ideal is never far from principal.

\begin{proposition}[one and a half generators]\label{prop:twogen}
Let $R$ be Dedekind, $I$ a non-zero ideal and $0\ne a\in I$. Then $I=(a,b)$ for some $b\in I$.
\end{proposition}

\begin{proof}
Let $\mathfrak p_1,\dots,\mathfrak p_r$ be the primes containing $a$, those with $v_{\mathfrak p}(aR)>0$. Every prime dividing $I$ contains $a$, so $I=\prod_i\mathfrak p_i^{n_i}$ with $n_i=v_{\mathfrak p_i}(I)$. Since $\mathfrak p_i^{n_i+1}\ne\mathfrak p_i^{n_i}$, we
can choose $b_i\in\mathfrak p_i^{n_i}\setminus\mathfrak p_i^{n_i+1}$. The ideals $\mathfrak p_i^{n_i+1}$ are pairwise coprime (no maximal ideal contains two of them), so the Chinese remainder theorem (Volume~3, Theorem~6.9, and induction) gives
$b\in R$ with $b\equiv b_i\pmod{\mathfrak p_i^{n_i+1}}$ for every $i$. Then $b\in\mathfrak p_i^{n_i}\setminus\mathfrak p_i^{n_i+1}$, that is $v_{\mathfrak p_i}(bR)=n_i$, for each $i$; in particular $b\in\bigcap_i\mathfrak p_i^{n_i}=I$ by
Corollary~\ref{cor:contain}. By the same corollary, $v_{\mathfrak p}(aR+bR)=\min(v_{\mathfrak p}(aR),v_{\mathfrak p}(bR))$ equals $n_i$ at $\mathfrak p=\mathfrak p_i$ (as $v_{\mathfrak p_i}(aR)\ge n_i$, because $a\in I$), and $0$ at every other $\mathfrak p$. So
$aR+bR=I$.
\end{proof}

\section{Integral closure in a finite extension}

Here is the machine that produces Dedekind domains: start from one, extend its fraction field, and take integral elements.

\begin{lemma}\label{lem:trace}
Let $L/K$ be a finite extension of degree $n$ and $y\in L$, with minimal polynomial $f$ of degree $d$ over $K$. The characteristic polynomial $\chi_y$ of the $K$-linear map $m_y\colon L\to L$, $z\mapsto yz$, is $f^{n/d}$. Its trace
is the \emph{trace} $\operatorname{Tr}_{L/K}(y)$. If $R$ is integrally closed with fraction field $K$ and $y$ is integral over $R$, then $\chi_y\in R[x]$ and $\operatorname{Tr}_{L/K}(y)\in R$.
\end{lemma}

\begin{proof}
Let $\ell_1,\dots,\ell_k$ be a basis of $L$ over $K(y)$, so $k=n/d$ by the tower law (Volume~3, Theorem~8.2). In the $K$-basis $(y^i\ell_j)_{0\le i<d,\,1\le j\le k}$, the matrix of $m_y$ is block diagonal with $k$ copies of the
companion matrix of $f$, whose characteristic polynomial is $f$. For the second claim, $y$ is a root of a monic $g\in R[x]$, and $f\mid g$ in $K[x]$. In a splitting field of $g$, the roots of $f$ are roots of $g$, hence integral
over $R$; the coefficients of $f$ are, up to sign, elementary symmetric functions of these roots, so they are integral over $R$ (Corollary~\ref{cor:integralclosure}) and lie in $K$, hence in $R$.
\end{proof}

\begin{theorem}\label{thm:intclosure}
Let $R$ be a Dedekind domain with fraction field $K$, $L/K$ a finite separable extension of degree $n$, and $S$ the integral closure of $R$ in $L$. Then:
\begin{enumerate}[label=(\alph*)]
\item there is $\theta\in S$ with $L=K(\theta)$, and $dS\subseteq R[\theta]$ where $d=\det\bigl(\operatorname{Tr}_{L/K}(\theta^{i+j})\bigr)_{0\le i,j<n}$ is non-zero;
\item $S$ is a finitely generated $R$-module, and a Dedekind domain with fraction field $L$.
\end{enumerate}
\end{theorem}

\begin{proof}
(a) A Dedekind domain is infinite (a finite domain is a field), so $K$ is infinite. For a separable extension of an infinite field, the proof of Volume~3, Theorem~8.19 (the primitive element theorem) applies word for word,
choosing the constant $c$ in $K$: it only uses that the conjugates of an element are distinct. So $L=K(\theta)$. If $\theta$ has minimal polynomial $\sum_ic_ix^i$ and $r\in R$ is a common denominator of the $c_i$, then
$r\theta$ is a root of a monic polynomial over $R$ and still generates $L$; so we may take $\theta\in S$.

Let $\theta_1,\dots,\theta_n$ be the roots of the minimal polynomial of $\theta$ in a splitting field, distinct by separability. The matrix of $m_\theta$ in the basis $1,\theta,\dots,\theta^{n-1}$ is a companion matrix with distinct
eigenvalues $\theta_k$, so it is diagonalisable over the splitting field, and $m_{\theta^m}=m_\theta^m$ has eigenvalues $\theta_k^m$. Hence $\operatorname{Tr}(\theta^m)=\sum_k\theta_k^m$, and the matrix
$\bigl(\operatorname{Tr}(\theta^{i+j})\bigr)$ is $V^{\mathsf T}V$ with the Vandermonde matrix $V=(\theta_k^i)$. So $d=\det(V)^2=\prod_{k<l}(\theta_l-\theta_k)^2\ne0$.

Now let $s\in S$ and write $s=\sum_jc_j\theta^j$ with $c_j\in K$. For each $i$, $\theta^is\in S$, so by Lemma~\ref{lem:trace}
\[
\operatorname{Tr}(\theta^is)=\sum_jc_j\operatorname{Tr}(\theta^{i+j})\in R .
\]
This is a linear system for the $c_j$ with matrix entries in $R$ and determinant $d$; by Cramer's rule $dc_j\in R$. So $ds\in R[\theta]$.

(b) By (a), $S\subseteq d^{-1}R[\theta]=\bigoplus_jRd^{-1}\theta^j$, a finitely generated $R$-module; since $R$ is Noetherian, its submodule $S$ is finitely generated (Proposition~\ref{prop:noethmodules}). The ideals of $S$ are
$R$-submodules, so $S$ is a Noetherian ring. Every $y\in L$ has $ry\in S$ for some $0\ne r\in R$ (as for $\theta$ above), so $L$ is the fraction field of $S$. If $y\in L$ is integral over $S$, it is integral over $R$ by
Corollary~\ref{cor:integralclosure}, so $y\in S$: $S$ is integrally closed. Finally $\dim S=\dim R=1$ by Corollary~\ref{cor:dimintegral}, so $S$ is not a field and a non-zero prime of $S$ cannot lie in a larger prime, since
$0\subsetneq\mathfrak q\subsetneq\mathfrak q'$ would be a chain of length $2$.
\end{proof}

With $R=\Z$ this gives once more that $\mathcal O_K$ is Dedekind; with $R=\C[x]$ and $L=\C(x)[y]/(y^2-x^3+x)$ it shows that the ring of functions on the cubic curve $y^2=x^3-x$ is Dedekind, provided one knows that this
ring, $\C[x,y]/(y^2-x^3+x)$, is the integral closure of $\C[x]$ in $L$ (it is, because the curve has no singular point; see Exercise~\ref{ex:cubicnotprincipal}). The separability hypothesis is not decorative: in characteristic $p$, the
integral closure of a Dedekind domain in an inseparable extension is still Dedekind (the Krull--Akizuki theorem), but it need not be a finitely generated $R$-module.

\section{How primes split}

Let $R$, $K$, $L$, $S$ be as in Theorem~\ref{thm:intclosure} and $\mathfrak p\ne0$ a prime of $R$. In $S$,
\[
\mathfrak pS=\mathfrak Q_1^{e_1}\cdots\mathfrak Q_g^{e_g}
\]
with distinct primes $\mathfrak Q_i$, the primes of $S$ \emph{above} $\mathfrak p$; indeed $\mathfrak Q_i\cap R$ is a prime containing $\mathfrak p$ and not containing $1$, so it equals $\mathfrak p$. The \emph{ramification index} is $e_i$, and
the \emph{residue degree} $f_i$ is the dimension of the field $S/\mathfrak Q_i$ over $R/\mathfrak p$.

\begin{theorem}[the fundamental identity]\label{thm:fundid}
$\displaystyle\sum_{i=1}^ge_if_i=n$.
\end{theorem}

\begin{proof}
\emph{The left side is $\dim_{R/\mathfrak p}S/\mathfrak pS$.} By the Chinese remainder theorem, $S/\mathfrak pS\cong\prod_iS/\mathfrak Q_i^{e_i}$, the $\mathfrak Q_i^{e_i}$ being pairwise coprime. In the chain
$S\supsetneq\mathfrak Q\supsetneq\mathfrak Q^2\supsetneq\cdots\supsetneq\mathfrak Q^e$, each quotient $\mathfrak Q^k/\mathfrak Q^{k+1}$ is a vector space over $S/\mathfrak Q$, and it has dimension $1$: for $x\in\mathfrak Q^k\setminus\mathfrak Q^{k+1}$,
the ideal $\mathfrak Q^{k+1}+xS$ lies in $\mathfrak Q^k$ and strictly contains $\mathfrak Q^{k+1}$ (the powers are distinct by unique factorisation); by Corollary~\ref{cor:contain} there is no ideal strictly between, so it is $\mathfrak Q^k$, and $x$ spans $\mathfrak Q^k/\mathfrak Q^{k+1}$. So $\dim_{R/\mathfrak p}S/\mathfrak Q^e=ef$.

\emph{The right side.} Localise at $\mathfrak p$: $R_{\mathfrak p}$ is a DVR, and $S_{\mathfrak p}=(R\setminus\mathfrak p)^{-1}S$ is the integral closure of $R_{\mathfrak p}$ in $L$ (if $y$ is integral over $R_{\mathfrak p}$, a common denominator
$t\notin\mathfrak p$ of its equation makes $ty$ integral over $R$). By Theorem~\ref{thm:intclosure}(b) applied to the Dedekind domain $R_{\mathfrak p}$, $S_{\mathfrak p}$ is a finitely generated, torsion-free module over the
principal ideal domain $R_{\mathfrak p}$, hence free (Volume~3, Theorem~7.6). A basis of $S_{\mathfrak p}$ over $R_{\mathfrak p}$ is linearly independent over $K$ (clear denominators) and spans $L$ (every $y\in L$ has a multiple $ry\in S$),
so it has $n$ elements, and
\[
S_{\mathfrak p}/\mathfrak pS_{\mathfrak p}\cong(R_{\mathfrak p}/\mathfrak pR_{\mathfrak p})^n\cong(R/\mathfrak p)^n.
\]
Finally, the natural map $S/\mathfrak pS\to S_{\mathfrak p}/\mathfrak pS_{\mathfrak p}$ is an isomorphism. For $t\in R\setminus\mathfrak p$, pick $u\in R$ with $tu\equiv1\pmod{\mathfrak p}$ ($R/\mathfrak p$ is a field). Then $s/t\equiv us$
modulo $\mathfrak pS_{\mathfrak p}$, which gives surjectivity; and if $s\in\mathfrak pS_{\mathfrak p}$, then $ts\in\mathfrak pS$ for some such $t$, and $s=u(ts)-(tu-1)s\in\mathfrak pS$, which gives injectivity.
\end{proof}

For $R=\Z$ this is Volume~13, Theorem~4.11, proved there with norms of ideals; here it comes from linear algebra over a DVR. When moreover $L/K$ is Galois, the Galois group permutes the $\mathfrak Q_i$ transitively, all $e_i$ are
equal and all $f_i$ are equal, so $efg=n$: Volume~13, Chapter~4, proves this for number fields, and the proof given there works verbatim over any Dedekind domain.

\section{Numbers and curves}

\begin{figure}[ht]
\centering
\fig{ch05-cubicclass}
\caption{The real points of the cubic $E\colon y^2=x^3-x+1$. The line $y=1-2x$ meets $E$ at $P=(0,1)$, $Q=(1,-1)$ and $R=(3,-5)$, so the ideal it generates in $\C[x,y]/(y^2-x^3+x-1)$ is $\mathfrak m_P\mathfrak m_Q\mathfrak m_R$.
The vertical line $x=3$ gives $(x-3)=\mathfrak m_R\mathfrak m_{R'}$ with $R'=(3,5)$. In the class group, $[\mathfrak m_P][\mathfrak m_Q]=[\mathfrak m_{R'}]$: the chord-and-tangent addition of points.}
\label{fig:cubicclass}
\end{figure}

Dedekind and Weber's insight was that the ring $A=\C[x,y]/(y^2-g(x))$ of a curve, with $g$ a cubic with distinct roots, behaves like a ring of integers. Its maximal ideals are the $\mathfrak m_P=(x-a,y-b)$ for the points
$P=(a,b)$ of the curve (Nullstellensatz, Theorem~\ref{thm:nullstellensatz}). A line $\ell=y-\lambda x-\mu$ meets the curve where $(\lambda x+\mu)^2=g(x)$, at three points $P,Q,R$, and if they are distinct the factorisation of
$\ell A$ is $\mathfrak m_P\mathfrak m_Q\mathfrak m_R$: each $\mathfrak m$ containing $\ell$ is one of the three, and $A/\ell A\cong\C[x]/\bigl((\lambda x+\mu)^2-g(x)\bigr)$ has dimension $3$ over $\C$, which, as in the proof of
Theorem~\ref{thm:fundid}, is the sum of the exponents. A vertical line gives $(x-a)=\mathfrak m_P\mathfrak m_{P'}$, where $P'=(a,-b)$ is the mirror image of $P=(a,b)$.

So in $\operatorname{Cl}(A)$, $[\mathfrak m_P][\mathfrak m_Q][\mathfrak m_R]=1$ and $[\mathfrak m_P][\mathfrak m_{P'}]=1$. These are exactly the rules of the chord-and-tangent addition on an elliptic curve (Volume~13, where the group law is proved): the
third point of the line through $P$ and $Q$, reflected, is $P+Q$. The factorisations show that $P\mapsto[\mathfrak m_P]$, with $O\mapsto1$ for the point at infinity, respects those rules when the points of the line are distinct. They do not show that the map is bijective, so it is not claimed to be an isomorphism onto $\operatorname{Cl}(A)$. What the calculation does show is a non-principal ideal: a point of order $2$ such as $(0,0)$ on $y^2=x^3-x$ gives an
ideal $\mathfrak m=(x,y)$ that is not principal although $\mathfrak m^2=(x)$ is (Exercise~\ref{ex:cubicnotprincipal}).

\begin{beauty}
On 26 March 1940, imprisoned in Rouen for failing to report for military service, André Weil wrote to his sister, the philosopher Simone Weil, about mathematics. He described three ``languages'', number fields, function fields
over finite fields, and Riemann surfaces, as a Rosetta stone: a theorem in one column suggests, and sometimes proves, a theorem in another. This chapter is a grammar common to all three columns: Dedekind domains, prime
ideals, ramification and class groups mean the same thing for $\Z[\sqrt{-5}]$, for $\mathbb F_p[x]$ and its extensions, and for the complex cubic of Figure~\ref{fig:cubicclass}, a Riemann surface. Weil's own proof of the Riemann hypothesis for curves over finite fields, which he announced in 1940,
lives in the second column; the first is still waiting.
\end{beauty}

\begin{think}
The ring of the line, $\C[x]$, is a principal ideal domain, and its class group is trivial; the ring of the cubic has class group a torus. What would you guess for the ring of the circle $\R[x,y]/(x^2+y^2-1)$? (It is the ring
of trigonometric polynomials; Exercise~\ref{ex:circle} computes its class group.)
\end{think}

\section{Exercises}

\begin{exercise}[$\star$]\label{ex:zsqrt5}
In $R=\Z[\sqrt{-5}]$, with the norm $N(a+b\sqrt{-5})=a^2+5b^2$, prove the four factorisations of Example~\ref{ex:sixtwice}, and show that $\mathfrak p=(2,1+\sqrt{-5})$ is not principal.
\end{exercise}

\begin{exercise}[$\star$]\label{ex:powerseries}
Let $k$ be a field. Show that $k[[t]]$, the ring of formal power series, is a DVR, and describe its valuation.
\end{exercise}

\begin{exercise}[$\star\star$]\label{ex:semilocal}
Show that a Dedekind domain with only finitely many prime ideals is a principal ideal domain.
\end{exercise}

\begin{exercise}[$\star\star$]\label{ex:notdedekind}
Show that $\C[x,y]$ and $\Z[\sqrt{-3}]$ are Noetherian domains that are not Dedekind, and in each case find a non-zero prime ideal that is not invertible.
\end{exercise}

\begin{exercise}[$\star\star$]\label{ex:circle}
Let $A=\R[x,y]/(x^2+y^2-1)$, a Dedekind domain (take this for granted). Show that $\mathfrak m=(x-1,y)$ is not principal but $\mathfrak m^2=(x-1)$, and that $A\otimes_\R\C\cong\C[u,u^{-1}]$ is a principal ideal domain.
(In fact $\operatorname{Cl}(A)\cong\Z/2$.)
\end{exercise}

\begin{exercise}[$\star\star\star$]\label{ex:cubicnotprincipal}
Let $A=\C[x,y]/(y^2-x^3+x)$, and $N(a+by)=a^2-b^2(x^3-x)\in\C[x]$ for $a,b\in\C[x]$.
\begin{enumerate}[label=(\alph*)]
\item Show that every element of $A$ is uniquely $a+by$ with $a,b\in\C[x]$, that $N$ is multiplicative, and that the units of $A$ are the non-zero constants.
\item Show that $\mathfrak m=(x,y)$ satisfies $\mathfrak m^2=(x)$, and that it is not principal. (Look at the degree of $N$.)
\item Show that $A$ is integrally closed.
\end{enumerate}
\end{exercise}

\section{Hints}

\begin{hint}{ex:zsqrt5}
Compute the products of generators, and look for a generator of the result. A generator $g$ of $\mathfrak p$ would have $N(g)$ dividing $N(2)$ and $N(1+\sqrt{-5})$.
\end{hint}

\begin{hint}{ex:powerseries}
A series with non-zero constant term is invertible.
\end{hint}

\begin{hint}{ex:semilocal}
By the Chinese remainder theorem, find $\pi$ with $v_{\mathfrak p_1}(\pi)=1$ and $v_{\mathfrak p_j}(\pi)=0$ for $j\ne1$.
\end{hint}

\begin{hint}{ex:notdedekind}
For $\Z[\sqrt{-3}]$, think of $\frac{1+\sqrt{-3}}2$, and compute $\mathfrak p^2$ for $\mathfrak p=(2,1+\sqrt{-3})$.
\end{hint}

\begin{hint}{ex:circle}
Use the norm $N(a+by)=a^2+b^2(x^2-1)\in\R[x]$ and look at the parity of its degree. Over $\C$, put $u=x+iy$.
\end{hint}

\begin{hint}{ex:cubicnotprincipal}
In $N(a+by)$, the terms $a^2$ and $b^2(x^3-x)$ have degrees of different parity. For (c), an element $\alpha$ of $\C(x)[y]$ integral over $A$ is integral over $\C[x]$, so its trace $2a$ and norm $a^2-b^2(x^3-x)$ lie in $\C[x]$.
\end{hint}

\section{Solutions}

\begin{solution}{ex:zsqrt5}
\idea Multiply generators, then simplify by subtracting; for non-principality, use the norm, which is multiplicative and takes only certain values.
\writeup $\mathfrak p^2$ is generated by $4$, $2(1+\sqrt{-5})$ and $(1+\sqrt{-5})^2=-4+2\sqrt{-5}$. All are multiples of $2$, and $2+2\sqrt{-5}-(-4+2\sqrt{-5})=6$, so $6-4=2\in\mathfrak p^2$: hence $\mathfrak p^2=2R$.
Similarly $\mathfrak q\bar{\mathfrak q}$ is generated by $9$, $3(1-\sqrt{-5})$, $3(1+\sqrt{-5})$ and $(1+\sqrt{-5})(1-\sqrt{-5})=6$, all multiples of $3$, and $9-6=3$ is among their combinations; so $\mathfrak q\bar{\mathfrak q}=3R$.
$\mathfrak p\mathfrak q$ is generated by $6=(1+\sqrt{-5})(1-\sqrt{-5})$, $2(1+\sqrt{-5})$, $3(1+\sqrt{-5})$ and $(1+\sqrt{-5})^2$, all multiples of $1+\sqrt{-5}$, and $3(1+\sqrt{-5})-2(1+\sqrt{-5})=1+\sqrt{-5}$; so $\mathfrak p\mathfrak q=(1+\sqrt{-5})$.
Conjugating gives $\mathfrak p\bar{\mathfrak q}=(1-\sqrt{-5})$, since $\bar{\mathfrak p}=(2,1-\sqrt{-5})=\mathfrak p$ (because $1-\sqrt{-5}=2-(1+\sqrt{-5})$).

If $\mathfrak p=gR$, then $g$ divides $2$ and $1+\sqrt{-5}$, so $N(g)$ divides $N(2)=4$ and $N(1+\sqrt{-5})=6$, hence $N(g)\in\{1,2\}$. But $a^2+5b^2=2$ has no integer solution, and $N(g)=1$ means $g=\pm1$ and $\mathfrak p=R$,
which is false: in $R/\mathfrak p$, $\sqrt{-5}\equiv-1$, so every element is congruent to an integer, and $2\equiv0$, so $R/\mathfrak p\cong\mathbb F_2\ne0$.
\end{solution}

\begin{solution}{ex:powerseries}
\idea The order of vanishing at $t=0$ is the valuation.
\writeup If $f=\sum a_it^i$ with $a_0\ne0$, solve $fg=1$ for $g=\sum b_it^i$ recursively: $b_0=a_0^{-1}$ and $b_n=-a_0^{-1}\sum_{i=1}^na_ib_{n-i}$. A non-zero $f$ is $t^nu$, where $n$ is the index of its first non-zero coefficient and
$u$ is a unit. The product of two non-zero series has first coefficient the product of the first coefficients, so $k[[t]]$ is a domain. A non-zero ideal is $(t^n)$ with $n$ the least order of its elements, so $k[[t]]$ is a principal ideal
domain whose only non-zero prime is $(t)$. The valuation is $v(f)=$ the order of $f$, extended to the field of formal Laurent series $k((t))$.
\end{solution}

\begin{solution}{ex:semilocal}
\idea A prime that is the only one in sight can be generated by a uniformiser that avoids the others.
\writeup Let $\mathfrak p_1,\dots,\mathfrak p_r$ be the non-zero primes. Choose $\pi_0\in\mathfrak p_1\setminus\mathfrak p_1^2$. By the Chinese remainder theorem (the ideals $\mathfrak p_1^2,\mathfrak p_2,\dots,\mathfrak p_r$ are pairwise coprime) there
is $\pi$ with $\pi\equiv\pi_0\pmod{\mathfrak p_1^2}$ and $\pi\equiv1\pmod{\mathfrak p_j}$ for $j\ge2$. Then $v_{\mathfrak p_1}(\pi R)=1$ and $v_{\mathfrak p_j}(\pi R)=0$, so $\pi R=\mathfrak p_1$. The same holds for every
$\mathfrak p_i$, so all primes are principal, hence all ideals, which are products of primes.
\end{solution}

\begin{solution}{ex:notdedekind}
\idea One ring has dimension $2$; the other is not integrally closed.
\writeup $\C[x,y]$ is Noetherian (Theorem~\ref{thm:hilbertbasis}), and $(x)\subsetneq(x,y)$ is a non-zero prime that is not maximal. The prime $(x,y)$ is not invertible: if it were, its localisation at itself would be
principal (Lemma~\ref{lem:invertible}), so by Nakayama's lemma (Corollary~\ref{cor:mingens}) the space $\mathfrak m/\mathfrak m^2$, which has basis $x,y$ over $\C$, would have dimension at most $1$.

In $R=\Z[\sqrt{-3}]$, $\omega=\frac{1+\sqrt{-3}}2$ is a root of $x^2-x+1$ and lies in the fraction field but not in $R$, so $R$ is not integrally closed. Let $\mathfrak p=(2,1+\sqrt{-3})$. In $R/\mathfrak p$, $\sqrt{-3}\equiv-1$
and $2\equiv0$, so $R/\mathfrak p\cong\mathbb F_2$ and $\mathfrak p$ is maximal. Now $\mathfrak p^2$ is generated by $4$, $2+2\sqrt{-3}$ and $(1+\sqrt{-3})^2=-2+2\sqrt{-3}=2(-1+\sqrt{-3})$, so $\mathfrak p^2=2\mathfrak p$, as
$-1+\sqrt{-3}=(1+\sqrt{-3})-2\in\mathfrak p$ and conversely $2,1+\sqrt{-3}$ are combinations of $2,-1+\sqrt{-3}$. If $\mathfrak p$ were invertible, multiplying by $\mathfrak p^{-1}$ would give $\mathfrak p=2R$, but
$1+\sqrt{-3}\notin2R$.
\end{solution}

\begin{solution}{ex:circle}
\idea As for the cubic, the norm $N(a+by)=a^2-b^2(1-x^2)$ is a polynomial in $x$, and over $\R$ it cannot have degree $1$.
\writeup Every element of $A$ is uniquely $a+by$ with $a,b\in\R[x]$, since $y^2=1-x^2$, and $y\mapsto-y$ is an automorphism, so $N(\alpha)=\alpha\bar\alpha$ is multiplicative. We have $\mathfrak m^2=((x-1)^2,(x-1)y,y^2)$
and $y^2=-(x-1)(x+1)$, so $\mathfrak m^2\subseteq(x-1)$; conversely $x-1=-\tfrac12\bigl((x-1)^2+y^2\bigr)\in\mathfrak m^2$ (as $(x-1)^2+y^2=x^2-2x+1+1-x^2=-2(x-1)$). So $\mathfrak m^2=(x-1)$.

If $\mathfrak m=(g)$, then $N(g)$ divides $N(x-1)=(x-1)^2$ and $N(y)=x^2-1$ in $\R[x]$, so $N(g)$ is a constant or a constant times $x-1$; it is not a constant, since $g$ would then be a unit and $\mathfrak m=A$, while
$A/\mathfrak m\cong\R$. But $N(a+by)=a^2+b^2(x^2-1)$ has leading term $\alpha^2x^{2p}$ or $\beta^2x^{2q+2}$ or their sum, a positive multiple of an even power of $x$ (the two leading coefficients, when the degrees agree, are both
positive, so they cannot cancel). So $N(g)$ has even degree, a contradiction. Over $\C$, with $u=x+iy$, one has $u\cdot(x-iy)=x^2+y^2=1$, so $A\otimes\C=\C[u,u^{-1}]$, with $x=\frac{u+u^{-1}}2$, $y=\frac{u-u^{-1}}{2i}$; this is a
localisation of the principal ideal domain $\C[u]$, hence principal. The point is that over $\R$ the circle has no real points at infinity, but over $\C$ it has two, $u=0$ and $u=\infty$, and they kill the class of $\mathfrak m$.
\end{solution}

\begin{solution}{ex:cubicnotprincipal}
\idea The norm to $\C[x]$ turns questions about $A$ into questions about degrees of polynomials, and degrees of norms are never $1$.
\writeup (a) Put $h=x^3-x$. As a $\C[x]$-module, $\C[x][Y]/(Y^2-h)$ is free on $1,Y$ (division by the monic $Y^2-h$), and this is $A$. The automorphism $y\mapsto-y$ gives $N(\alpha)=\alpha\bar\alpha$, so $N$ is multiplicative.
If $\alpha$ is a unit, $N(\alpha)$ is a unit of $\C[x]$, a non-zero constant. The degree of $a^2$ is even and that of $b^2h$ is odd (when $b\ne0$), so they cannot cancel and $\deg N(a+by)=\max(2\deg a,2\deg b+3)$; for this to be
$0$ we need $b=0$ and $a$ constant.

(b) $\mathfrak m^2=(x^2,xy,y^2)$ and $y^2=x(x^2-1)$, so $\mathfrak m^2\subseteq(x)$, and $x=-y^2+x\cdot x^2\in\mathfrak m^2$ since $-y^2+x^3=x$. Hence $\mathfrak m^2=(x)$. If $\mathfrak m=(g)$, then $N(g)$ divides
$N(x)=x^2$ and $N(y)=-h=-x(x-1)(x+1)$, so $N(g)$ divides $x$; it is not constant, as $g$ is not a unit ($A/\mathfrak m\cong\C$). So $\deg N(g)=1$, impossible by (a).

(c) Let $\alpha=a+by$ with $a,b\in\C(x)$ be integral over $A$; then it is integral over $\C[x]$ ($A$ is integral over $\C[x]$, Corollary~\ref{cor:integralclosure}). Its conjugate $\bar\alpha$ is integral too, so $2a=\alpha+\bar\alpha$ and
$N(\alpha)=a^2-b^2h$ lie in $\C(x)$ and are integral over the integrally closed $\C[x]$ (a principal ideal domain), so they are in $\C[x]$. Thus $a\in\C[x]$ and $b^2h\in\C[x]$. Since $h$ is square-free, a prime factor of
the denominator of $b$ would appear squared in the denominator of $b^2$ and could not be cancelled by $h$; so $b\in\C[x]$ and $\alpha\in A$. Square-freeness of $h$ is exactly the absence of singular points: for $y^2=x^3$ the
element $y/x$ is integral ($(y/x)^2=x$) but not in the ring.
\end{solution}

\section*{Sources and further reading}
\addcontentsline{toc}{section}{Sources and further reading}

R.~Dedekind, Supplement~X to P.~G.~Lejeune Dirichlet, \emph{Vorlesungen über Zahlentheorie}, 2nd edition (Vieweg, 1871), rewritten in the 3rd (1879) and 4th (1894) editions; R.~Dedekind and H.~Weber, ``Theorie der algebraischen
Functionen einer Veränderlichen'', \emph{J.\ Reine Angew.\ Math.} 92 (1882), 181--290; E.~Noether, ``Abstrakter Aufbau der Idealtheorie in algebraischen Zahl- und Funktionenkörpern'', \emph{Math.\ Ann.} 96 (1927), 26--61;
L.~Claborn, ``Every abelian group is a class group'', \emph{Pacific J.\ Math.} 18 (1966), 219--222. The story of the calendar is told by E.~T.~Bell, \emph{Men of Mathematics} (1937), in the chapter on Kummer and Dedekind;
Noether's remark is discussed by L.~Corry, ``Steht es alles wirklich schon bei Dedekind? Ideals and factorization between Dedekind and Noether'' (2017). Weil's letter is translated by M.~Krieger, ``A 1940 letter of André Weil on
analogy in mathematics'', \emph{Notices Amer.\ Math.\ Soc.} 52 (2005), 334--341.

For more: M.~F.~Atiyah and I.~G.~Macdonald, \emph{Introduction to Commutative Algebra} (1969), Chapter~9; J.~Neukirch, \emph{Algebraic Number Theory} (Springer, 1999), Chapter~I, \S\S3 and 8; H.~Matsumura, \emph{Commutative Ring Theory}
(1986), Chapter~11, for Krull--Akizuki and the general theory; J.~H.~Silverman, \emph{The Arithmetic of Elliptic Curves} (2nd edition, 2009), Chapter~III, for the group law and the Picard group of a curve.
